% !PRINTED 199
\weberpage{199}

%<<CONT
\weq{12}{P_1 = 1,\quad P_2 = -1,\quad P_3 = 3\wp(u)^4 - \frac32 g_2\wp(u)^2 - 3g_3\wp(u) - \frac{g_2^2}{16}}
wird. Wir erhalten ferner aus der Additionsformel (§ 49):
\[
\wp(v+u) - \wp(v-u) = -\frac{\wp'(u)\,\wp'(v)}{\left[\wp(v)-\wp(u)\right]^2},
\]
indem wir $v=2u$ setzen, nach (9):
\[
\wp(3u) - \wp(u) = -\frac{\wp'(u)\,\wp'(2u)}{\left[\wp(2u)-\wp(u)\right]^2}
= -\frac{\wp'(u)^5\,\wp'(2u)}{\psi_3(u)^2} = -\frac{\psi_2(u)\,\psi_4(u)}{\psi_3(u)^2},
\]
also:
\[
\psi_4(u) = \wp'(u)\,P_4 = -\wp'(u)^4\,\wp'(2u),
\]
und wenn man aus (11) den Wert $\wp'(2u)$ bildet:
\weq{13}{P_4 = -2\wp(u)^6 + \frac52 g_2\wp(u)^4 + 10g_3\,p(u)^3 % [sic]
+ \frac58 g_2^2\wp(u)^2 + \frac12 g_2g_3\wp(u) + g_3^2 - \frac{g_2^3}{32}.}

Für größere Werte von $n$ leiten wir Rekursionsformeln ab. Zu diesem Zweck
wenden wir die Formel (9) auf zwei verschiedene Zahlen $m, n$ an und
erhalten:
\[
\wp(u) - \wp(mu) = \frac{\psi_{m+1}(u)\,\psi_{m-1}(u)}{\psi_m(u)^2},
\]
\[
\wp(u) - \wp(nu) = \frac{\psi_{n+1}(u)\,\psi_{n-1}(u)}{\psi_n(u)^2},
\]
woraus zu ersehen ist, daß die rechten Seiten einander gleich werden für
solche Werte $u^0$ von $u$, die von Null verschieden sind und der
Bedingung genügen
\[
m\,u^0 \equiv \pm n\,u^0 \pmod{\omega_1,\omega_2},
\]
oder
\[
(m\pm n)\,u^0 \equiv 0 \pmod{\omega_1,\omega_2};
\]
für dieselben Werte $u^0$ verschwinden aber auch die Funktionen
\[
\psi_{m\pm n}(u),
\]
so daß die beiden Funktionen:
\[
\psi_{m+1}(u)\,\psi_{m-1}(u)\,\psi_n(u)^2 - \psi_{n+1}(u)\,\psi_{n-1}(u)\,\psi_m(u)^2,
\]
\[
\psi_{m+n}(u)\,\psi_{m-n}(u),
\]
die nach (6) ganze rationale Funktionen von $\wp(u)$ sind, für die
nämlichen endlichen Werte von $\wp(u)$ verschwinden. Sie unter-
%CONT>>
